\scanpage{90}
% L1–2, continuation of brace
\[
\left\{\begin{aligned}
\Theta_3(v\mid t)&=\sum_{-\infty}^{+\infty}e^{n^2i\pi t+2ni\pi v}\\
&=1+2\sum_1^\infty e^{n^2i\pi t}\cos2n\pi v=\vartheta_3(v)
\end{aligned}\right.
\]

% L3
($\Theta(v\mid t)=\vartheta_1(v)$ \quad -- $\vartheta_i$, fonctions de Weierstrass)

% L4–6
Les autres quantités de la théorie des fonctions elliptiques se calculent avec \unsure{avantage|au moyen} à partir des fonctions $\Theta$, une fois fixés $\omega_1$ et $t=\frac{\omega_2}{\omega_1}$.

% L7
Le calcul des $e_i$ par exemple donne

% L8
\[
e_i=\left(\frac1{2\omega_1}\right)^2
\left(\frac13\frac{\Theta'''0}{\Theta'0}-\frac{\Theta_i'''0}{\Theta_i'0}\right)
\quad\text{ou aussi}
\]

% L9–11
\[
\left\{\begin{aligned}
e_1&=\frac13\left(\frac{\pi}{2\omega_1}\right)^2\left[(\Theta_2 0)^4+(\Theta_3 0)^4\right]\\
e_2&=\frac13\left(\frac{\pi}{2\omega_1}\right)^2\left[-(\Theta_1 0)^4-(\Theta_3 0)^4\right]\\
e_3&=\frac12\left(\frac{\pi}{2\omega_1}\right)^2\left[(\Theta_1 0)^4-(\Theta_{\unsure{2|3}}0)^4\right]
\end{aligned}\right.
\tag*{\textcircled{\scriptsize 15}}
\]

% L12
Enfin, on a les relations

% L13
\[
2\eta_1\omega_1=-\frac{\Theta'''0}{6\Theta'0}\quad;\qquad
\Theta'0=\pi\,\Theta_1 0\,\Theta_2 0\,\Theta_3 0
\]

% lower left
% archive sheet number: 89
